% Generated draft: requires mathematical and source review.
\documentclass[11pt]{article}
\usepackage[margin=25mm]{geometry}
\usepackage{fontspec}
\setmainfont{Latin Modern Roman}
\usepackage{amsmath,amssymb,mathtools,unicode-math}
\setmathfont{Latin Modern Math}
\usepackage{enumitem,needspace,xurl,hyperref}
\setlength{\parindent}{0pt}
\setlength{\parskip}{5pt}
\setlength{\emergencystretch}{3em}
\newcommand{\N}{\mathbb{N}}
\newcommand{\Z}{\mathbb{Z}}
\newcommand{\Q}{\mathbb{Q}}
\newcommand{\R}{\mathbb{R}}
\newcommand{\C}{\mathbb{C}}
\newcommand{\F}{\mathbb{F}}
\newcommand{\A}{\mathbb{A}}
\newcommand{\PP}{\mathbb{P}}
\newcommand{\E}{\mathbb{E}}
\newcommand{\eps}{\varepsilon}
\newcommand{\dd}{\,\mathrm{d}}
\begin{document}
\section*{UnsolvedMath ID: 30005457}
\textbf{Source problem code:} catalogue entry
\subsection{Definitions and mathematical statement}
For an integer \(d\geq 2\), let \(\Z^d\) denote the set of integer \(d\)-tuples. If a reinforcement process on this vertex set determines a surviving-edge graph \(H=(\Z^d,E_{\mathrm{surv}})\), define
\[
\mathcal C_\infty(H)=\{C:C\text{ is a connected component of }H\text{ and }|V(C)|=\infty\}.
\]
Thus the property of having infinitely many infinite surviving components is
\[
|\mathcal C_\infty(H)|=\infty.
\]
This definition expresses the component property once the surviving-edge graph has been specified; it does not supply the missing definition of survival.

\textbf{Question (supplied Open Problem 3).} Can the regular-tree example referenced in the source be modified to give an example on \(\Z^d\) for some integer \(d\geq 2\)? The supplied expanded statement describes the intended example as a supercritical-reinforcement process with infinitely many infinite surviving components, so the requested modification is understood to retain those features.

\textbf{Unresolved formulation qualifications.} The supplied record does not define the reinforcement process, its parameters, the meaning of supercritical reinforcement, the initial conditions, the underlying lattice edges, or the criterion for an edge to survive. It also does not specify whether the displayed component property must hold almost surely, with positive probability, or under another probability qualification. Consequently, a fully quantified stochastic existence statement cannot be recovered from this record alone. The surviving-edge interpretation above is conditional on the intended meaning of survival and requires confirmation against the surrounding source discussion.
\subsection{Short English statement}
The supplied question asks whether the referenced regular-tree construction can be adapted to \(\Z^d\) for at least one integer \(d\geq 2\), retaining supercritical reinforcement and infinitely many infinite surviving components. The process, survival rule, and required probability qualification are missing, so the stochastic formulation remains incomplete.
\subsection{Sources}
\begin{itemize}
\item Local UnsolvedMath registry: \texttt{data/unsolvedmath/problems.json}, permanent ID \(30005457\), title ``Reinforcement Counterexamples on Integer Lattices,'' source code \texttt{catalogue entry}. The supplied record identifies the extracted passage as Open Problem 3 and gives the proposed year as 2023; no individual proposer is supplied.
\item Supplied source citation: \emph{MATRIX-MFO Tandem Workshop: Stochastic Reinforcement Processes and Graphs} (2023), Oberwolfach Reports, pp.\ 639--679, DOI \texttt{10.4171/owr/2023/12}. Source URL and literature-source link labelled ``original Oberwolfach report'': \url{https://doi.org/10.4171/owr/2023/12}.
\item Supplied literature source: Hirsch--Holmes--Kleptsyn, \emph{WARM percolation on a regular tree}, \url{https://arxiv.org/abs/2009.07682}.
\item Supplied literature assessment, checked on \texttt{2026-08-22}: the item concerns transferring the referenced regular-tree counterexample to \(\Z^d\) for some \(d\geq 2\); the curation reports finding neither such a construction nor an impossibility theorem and labels the problem open. This is supplied dated curation, not an independently verified literature assessment. The background explicitly calls for independent verification of time-sensitive status and cited claims.
\item The supplied registry labels the statement \texttt{corrected\_verified} and reports comparison with the official Oberwolfach report. These are supplied verification claims; no independent source consultation or mathematical verification is asserted here.
\end{itemize}
\end{document}
